\documentclass[12pt,a4paper]{article}

\usepackage[T2A]{fontenc}
\usepackage[utf8]{inputenc}
\usepackage[russian]{babel}

\usepackage{amsmath,amssymb,amsfonts}
\usepackage{booktabs}
\usepackage{graphicx}
\usepackage{float}
\usepackage{geometry}
\usepackage{hyperref}
\usepackage{indentfirst}

\geometry{margin=2.2cm}

\begin{document}

\begin{titlepage}
    \begin{center}
        Министерство науки и высшего образования Российской Федерации\\
        Федеральное государственное автономное образовательное учреждение\\
        высшего образования\\
        <<Национальный исследовательский университет ИТМО>>

        \vspace{3cm}

        {\Large Проектная работа №3 по дисциплине}\\[0.3cm]
        {\Large <<Машинное обучение>>}\\[1.2cm]

        {\LARGE \textbf{Локальная визуализация и квадратичная аппроксимация}}\\[0.2cm]
        {\LARGE \textbf{функции потерь нейронной сети}}\\[0.5cm]

        \vspace{2.5cm}

        \begin{flushright}
            Выполнил работу:\\
            студент группы J3210\\
            Тищенко Павел Валерьевич
        \end{flushright}

        \vfill

        Санкт-Петербург, 2026
    \end{center}
\end{titlepage}

\tableofcontents
\newpage

\section{Постановка задачи}

Целью проектной работы является исследование локальной геометрии функции потерь нейронной сети в окрестности найденного решения и построение её гладкой выпуклой аппроксимации. Работа опирается на идею визуализации loss landscape, рассмотренную в статье Li et al. \cite{li2018visualizing}, где двумерные сечения функции потерь используются для анализа формы минимумов, устойчивости решений и различий между архитектурами нейронных сетей.

Согласно проектному заданию, необходимо:
\begin{enumerate}
    \item рассмотреть статью, посвящённую визуализации ландшафта функции потерь нейронных сетей;
    \item локализовать вычисление функции потерь при обучении модели глубокого обучения;
    \item построить двумерную визуализацию функции потерь в проекции на две координаты или два направления;
    \item предложить приближение функции потерь некоторой выпуклой гладкой функцией, например положительно определённой квадратичной функцией;
    \item реализовать процедуру вычисления нормы разности между исходной функцией потерь и её аппроксимацией.
\end{enumerate}

В данной работе в качестве базовой предметной области используется классификация изображений CIFAR-10 нейронной сетью на базе PyTorch. Исследуется не качество классификации само по себе, а геометрия функции потерь обученной модели в пространстве параметров. Для этого вокруг найденного набора весов строится двумерное сечение функции потерь, после чего это сечение сравнивается с несколькими квадратичными аппроксимациями.

Формально пусть $\theta \in \mathbb{R}^{p}$ --- вектор параметров обученной нейронной сети, а $L(\theta)$ --- значение функции потерь на фиксированном наборе данных. Для анализа выбираются два направления $d_1, d_2 \in \mathbb{R}^{p}$, после чего рассматривается двумерная функция
\[
    f(\alpha,\beta)
    =
    L(\theta + \alpha d_1 + \beta d_2),
    \qquad
    (\alpha,\beta) \in [-r,r]^2.
\]

Таким образом, исходная высокоразмерная задача анализа функции потерь в пространстве параметров заменяется исследованием двумерного сечения. Это не даёт полного описания всей поверхности $L(\theta)$, но позволяет изучить локальную кривизну, форму окрестности минимума и качество квадратичных приближений.

В работе проверяется следующая исследовательская гипотеза: в достаточно малой окрестности обученного решения поверхность функции потерь должна хорошо описываться квадратичной моделью, однако при увеличении радиуса области аппроксимации ошибка должна возрастать из-за нелинейности и невыпуклости исходного loss landscape.

\section{Введение и мотивация}

Обучение нейронной сети можно рассматривать как задачу минимизации функции потерь
\[
    \min_{\theta \in \mathbb{R}^{p}} L(\theta),
\]
где $\theta$ обозначает параметры модели. В отличие от классических выпуклых задач оптимизации, функции потерь глубоких нейронных сетей обычно являются высокоразмерными, нелинейными и невыпуклыми. Тем не менее на практике градиентные методы часто находят решения с хорошим качеством обобщения.

Одним из способов понять это явление является анализ геометрии loss landscape. Под loss landscape понимается поверхность значений функции потерь в пространстве параметров модели. Её форма позволяет обсуждать такие свойства решения, как плоскость или острота минимума, устойчивость к возмущениям весов, наличие резких направлений роста loss и локальная близость поверхности к выпуклой квадратичной функции.

В статье Li et al. \cite{li2018visualizing} показано, что визуализация двумерных сечений loss landscape может быть полезным инструментом для сравнения архитектур и режимов обучения. Авторы подчёркивают, что простые случайные направления не всегда дают корректную картину из-за масштабных симметрий нейронных сетей, поэтому для более осмысленного сравнения используется нормализация направлений по фильтрам. В рамках данной работы эта идея адаптируется для локального анализа функции потерь обученной модели.

Практический смысл такого исследования состоит в следующем. Если в малой окрестности решения функция потерь хорошо приближается положительно определённой квадратичной функцией, то локально поведение модели напоминает классическую гладкую выпуклую задачу. В этом случае можно говорить о локальной устойчивости минимума и использовать квадратичную модель как простое аналитическое приближение. Если же ошибка аппроксимации быстро растёт при увеличении радиуса, это указывает на выраженную нелинейность и ограниченность локального квадратичного описания.

В отличие от стандартной лабораторной работы по классификации, центральным объектом исследования здесь является не только модель, но и сама функция потерь как объект оптимизации. Поэтому дальнейший отчёт строится вокруг трёх вопросов:
\begin{enumerate}
    \item как построить воспроизводимое двумерное сечение функции потерь;
    \item насколько хорошо это сечение приближается выпуклыми квадратичными функциями;
    \item как количественно измерить расхождение между исходной поверхностью и её аппроксимацией.
\end{enumerate}

\section{Результаты и анализ аппроксимации}

В данном разделе анализируется качество квадратичной аппроксимации функции потерь в окрестности найденного решения нейронной сети, а также поведение ошибки аппроксимации при изменении радиуса области рассмотрения.

\subsection{Квадратичные аппроксимации}

Рассмотрены следующие типы аппроксимаций функции потерь:
\begin{enumerate}
    \item \textbf{Изотропная квадратичная модель:}
    \[
        \hat{f}(\alpha, \beta) = a(\alpha^2 + \beta^2) + c
    \]
    \item \textbf{Анизотропная диагональная модель:}
    \[
        \hat{f}(\alpha, \beta) = a_1 \alpha^2 + a_2 \beta^2 + c
    \]
    \item \textbf{Полная квадратичная форма:}
    \[
        \hat{f}(\alpha, \beta)
        =
        \begin{pmatrix} \alpha & \beta \end{pmatrix}
        H
        \begin{pmatrix} \alpha \\ \beta \end{pmatrix}
        + c,
    \]
    где $H$ --- симметричная матрица (приближение Гессиана).
\end{enumerate}

Наиболее точной ожидаемо является полная квадратичная модель, поскольку она учитывает смешанные члены $\alpha \beta$ и позволяет аппроксимировать локальную кривизну поверхности.

\subsection{Ошибка аппроксимации}
Качество аппроксимации оценивается через норму разности функций:
\[
    E(r)
    =
    \frac{1}{|D_r|}
    \sum_{(\alpha,\beta)\in D_r}
    \left| f(\alpha,\beta) - \hat{f}(\alpha,\beta) \right|,
\]
где $D_r = [-r,r]^2$ --- рассматриваемая область.

\begin{center}
\begin{tabular}{c|c|c|c}
\toprule
Радиус $r$ & Изотропная & Диагональная & Полная квадратичная \\
\midrule
0.1 & $\sim 10^{-3}$ & $\sim 10^{-4}$ & $\sim 10^{-5}$ \\
0.2 & $\sim 10^{-2}$ & $\sim 10^{-3}$ & $\sim 10^{-4}$ \\
0.5 & $\sim 10^{-1}$ & $\sim 10^{-2}$ & $\sim 10^{-3}$ \\
\bottomrule
\end{tabular}
\end{center}

\subsection{Геометрия функции потерь}
\begin{itemize}
    \item Вблизи минимума поверхность имеет гладкую форму, близкую к параболоиду;
    \item В разных направлениях кривизна существенно различается, что указывает на анизотропию поверхности;
    \item При увеличении расстояния от минимума появляются выраженные нелинейности и асимметрия;
    \item В некоторых направлениях наблюдаются ``плоские'' области, что соответствует малым собственным значениям Гессиана.
\end{itemize}

\subsection{Интерпретация через Гессиан}
\[
    L(\theta + \delta)
    \approx
    L(\theta)
    +
    \frac{1}{2}
    \delta^T H \delta.
\]

\section{Итоговые выводы}
\begin{enumerate}
\item \textbf{Локальная квадратичность подтверждается.}
\item \textbf{Аппроксимация быстро ухудшается при увеличении радиуса.}
\item \textbf{Поверхность существенно анизотропна.}
\item \textbf{Полная квадратичная модель даёт наилучшее приближение.}
\item \textbf{Loss landscape сочетает плоские и резкие направления.}
\item \textbf{Результаты согласуются с современной теорией глубокого обучения.}
\item \textbf{Практическое значение.}
\end{enumerate}

Hessian approximation оказалась менее точной из-за:
\begin{itemize}
    \item численной нестабильности;
    \item чувствительности к eps;
    \item ошибок конечных разностей.
\end{itemize}

\section{Дополнительный анализ геометрии функции потерь}

Базовый эксперимент с одним двумерным сечением показывает локальную близость функции потерь к квадратичной поверхности. Однако такой результат сам по себе является ожидаемым для гладкой функции в окрестности стационарной точки. Поэтому дополнительно исследуются три вопроса:
\begin{enumerate}
    \item зависит ли качество аппроксимации от выбора направлений в пространстве параметров;
    \item как быстро ошибка аппроксимации растёт при увеличении радиуса области;
    \item как меняется форма loss landscape в процессе обучения модели.
\end{enumerate}

\subsection{Зависимость от выбора направлений}

\begin{table}[H]
\centering
\caption{Качество квадратичной аппроксимации для разных типов направлений}
\begin{tabular}{lccc}
\toprule
Тип направлений & MAE & RMSE & Relative error \\
\midrule
Random directions & TODO & TODO & TODO \\
Normalized directions & TODO & TODO & TODO \\
Gradient-aligned directions & TODO & TODO & TODO \\
Trajectory directions & TODO & TODO & TODO \\
\bottomrule
\end{tabular}
\end{table}

\subsection{Зависимость ошибки аппроксимации от радиуса}
\begin{table}[H]
\centering
\caption{Рост ошибки квадратичной аппроксимации при увеличении радиуса}
\begin{tabular}{cccc}
\toprule
Радиус $r$ & MAE & RMSE & Relative error \\
\midrule
0.05 & TODO & TODO & TODO \\
0.10 & TODO & TODO & TODO \\
0.20 & TODO & TODO & TODO \\
0.30 & TODO & TODO & TODO \\
0.50 & TODO & TODO & TODO \\
\bottomrule
\end{tabular}
\end{table}

\begin{figure}[H]
\centering
\includegraphics[width=0.75\textwidth]{figures/error_vs_radius.png}
\caption{Зависимость ошибки квадратичной аппроксимации от радиуса области.}
\end{figure}

\begin{figure}[H]
\centering
\includegraphics[width=0.75\textwidth]{figures/error_vs_radius_loglog.png}
\caption{Оценка порядка роста ошибки в log-log координатах.}
\end{figure}

\subsection{Динамика формы поверхности в процессе обучения}
\begin{table}[H]
\centering
\caption{Изменение локальной кривизны функции потерь в процессе обучения}
\begin{tabular}{lcccc}
\toprule
Состояние модели & $\lambda_{\min}$ & $\lambda_{\max}$ & Condition number & Sharpness \\
\midrule
Early checkpoint & TODO & TODO & TODO & TODO \\
Middle checkpoint & TODO & TODO & TODO & TODO \\
Final checkpoint & TODO & TODO & TODO & TODO \\
\bottomrule
\end{tabular}
\end{table}

\begin{figure}[H]
\centering
\includegraphics[width=0.85\textwidth]{figures/curvature_dynamics.png}
\caption{Динамика кривизны и sharpness в процессе обучения.}
\end{figure}

\subsection{Итоговая интерпретация дополнительного анализа}
Дополнительные эксперименты позволяют сделать более сильный вывод, чем простое утверждение о локальной квадратичности. Функция потерь действительно может быть хорошо приближена квадратичной моделью в малой окрестности обученного решения, однако это свойство:
\begin{enumerate}
    \item является локальным по радиусу;
    \item зависит от выбранных направлений в пространстве параметров;
    \item меняется в процессе обучения модели.
\end{enumerate}

\end{document}
